\documentclass[12pt]{article}

\usepackage{bbm}
\usepackage{fancyhdr}
\usepackage{titling}
\usepackage{enumitem}
\usepackage{amsmath}
\usepackage{dsfont}
\usepackage{amsfonts}
\usepackage{amssymb}
\usepackage{amsthm}
\usepackage{geometry}
\usepackage{tikz}
\usepackage{algorithm}
\usepackage{algpseudocode}
\geometry{verbose,tmargin=1in,bmargin=1in,lmargin=1in,rmargin=1in}
\renewcommand{\vec}[1]{\mathbf{#1}}
\newcommand{\dec}{\text{dec}}
\newcommand{\divi}{\mathbin{\textnormal{div}}}
\title{Homework 1}
\author{Satvik Prasad}
\date{}
\newtheorem{lemma}{Lemma}[section]
\fancypagestyle{plain}{%
  \fancyhf{}
  \fancyhead[L]{\theauthor}
  \fancyhead[R]{\thedate}
}
\pagestyle{plain}
\makeatletter
\def\@maketitle{%
  \newpage
  \null
  \vskip 1em%
  \begin{center}%
    \let \footnote \thanks
    {\LARGE \@title \par}%
    \vskip 1em%
  \end{center}%
  \par
\vskip 1em}
\makeatother
\usepackage{cmbright}

\newcommand{\mat}[1]{\mathbf{#1}}

% ---------- Probability theory ----------
% Core operators (optional argument-free; use as \E{X}, \Var{X}, \Cov{X}{Y})
\DeclareMathOperator{\EOp}{\mathbb{E}}
\DeclareMathOperator{\VarOp}{Var}
\DeclareMathOperator{\CovOp}{Cov}
\DeclareMathOperator{\CorrOp}{Corr}

\renewcommand{\P}[1]{\mathbb{P}\left(#1\right)}        % \P{A}
\newcommand{\E}[1]{\EOp\left[#1\right]}                % \E{X}
\newcommand{\Var}[1]{\VarOp\left(#1\right)}            % \Var{X}
\newcommand{\Cov}[2]{\CovOp\left(#1,\,#2\right)}       % \Cov{X}{Y}
\newcommand{\Corr}[2]{\CorrOp\left(#1,\,#2\right)}     % \Corr{X}{Y}
\newcommand{\SD}[1]{\operatorname{SD}\left(#1\right)}  % \SD{X}

% Subscripted versions, e.g. \Esub{\theta}{X}, \Psub{\theta}{A}
\newcommand{\Esub}[2]{\EOp_{#1}\left[#2\right]}
\newcommand{\Psub}[2]{\mathbb{P}_{#1}\left(#2\right)}

% Conditioning, indicators, independence
\newcommand{\given}{\,\middle|\,}                      % \E{X \given Y}
\newcommand{\ind}[1]{\mathbbm{1}\left\{#1\right\}}     % \ind{X > 0}
\newcommand{\indep}{\perp\!\!\!\perp}                  % X \indep Y
\newcommand{\iid}{\overset{\text{iid}}{\sim}}          % X_i \iid \Normal(0,1)

% Convergence
\newcommand{\convp}{\xrightarrow{\;p\;}}               % in probability
\newcommand{\convd}{\xrightarrow{\;d\;}}               % in distribution
\newcommand{\convas}{\xrightarrow{\text{a.s.}}}        % almost surely
\newcommand{\convLp}[1]{\xrightarrow{L^{#1}}}          % in L^p

% Probability spaces
\newcommand{\F}{\mathcal{F}}                           % sigma-algebra
\newcommand{\B}{\mathcal{B}}                           % Borel sigma-algebra
\newcommand{\pspace}{(\Omega, \F, \mathbb{P})}

% Common distributions
\newcommand{\Normal}{\mathcal{N}}
\newcommand{\Bern}{\operatorname{Bernoulli}}
\newcommand{\Bin}{\operatorname{Binomial}}
\newcommand{\Pois}{\operatorname{Poisson}}
\newcommand{\Unif}{\operatorname{Uniform}}
\newcommand{\Exp}{\operatorname{Exp}}
\newcommand{\Geom}{\operatorname{Geometric}}
\newcommand{\Gam}{\operatorname{Gamma}}
\newcommand{\Beta}{\operatorname{Beta}}

% Number sets
\newcommand{\R}{\mathbb{R}}
\newcommand{\N}{\mathbb{N}}
\newcommand{\Z}{\mathbb{Z}}


\begin{document}
\subsection{Writing Proofs}
\begin{enumerate}
  \item Assume we have finite dimensional $V,$ and a subspace $U$ of
    $V.$ Then, consider an arbitrary linear map $T \in \mathcal L(U,
    W)$ to a vector space $W.$

    Let $\beta_U = \{u_1, u_2, \ldots,u_m\}$ be a basis for $u$.
    Extend this basis to a basis for $V,$ $\{u_1, \ldots, u_n\}$
    where $n \geq m.$ Then, define the linear map $S \in \mathcal
    L(V, W)$ such that
    $$
    Su_k =
    \begin{cases}
      Tu_k & 0 \leq k \leq m \\
      0_W & m < k \leq n
    \end{cases}.
    $$

    Then, since $\beta_U$ is a basis, any $u \in U$ can be written as
    $$
    \mu_1u_1 + \ldots + \mu_m u_m,
    $$
    meaning
    \begin{align*}
      S u = \sum_{k = 1}^{m}\mu_k S u_k &= \sum_{k = 1}^{m} \mu_k Tu_k \\
      &= T\sum_{k = 1}^{m} \mu_ku_k = Tu \\
    \end{align*}
    by the linearity of $S, T.$ \qed

  \item Assume there is some $T \in \mathcal L(\mathbb R^5)$ such
    that $\text{null}(T) = \text{range}(T).$ Then, by the fundamental
    theorem of linear maps,
    $$
    \dim \mathbb R^5 = \dim \text{null}(T) + \dim \text{range}(T),
    $$
    or equivalently,
    $$
    5 = 2\dim \text{null}(T),
    $$
    meaning $5$ is even. However, this is a contradiction, and hence
    $\null(T) \neq \text{range}(T)$ for all linear maps $T \in
    \mathcal L(\mathbb R^5).$ \qed

  \item Consider an arbitrary polynomial $p \in \mathcal P(\mathbb
    R).$ Such a polynomial will have finite degree $n,$ meaning $p
    \in \mathcal P_n(\mathbb R).$ Then, consider the map $T \in
    \mathcal L(\mathcal P_{n + 1}(\mathbb R), \mathcal P_{n}(\mathbb
    R))$ defined such that
    $$
    Tq = 5q'' + 3q'.
    $$
    Consider the null-space of $T.$ Since $q$ is a degree-$n + 1$
    polynomial, we have
    $$
    q = \sum_{k = 0}^{n + 1} q_k x^k,
    $$
    meaning
    $$
    5q'' = \sum_{k = 0}^{n - 1} 5(k+2)(k + 1)q_{k + 2} x^{k}, \ 3q' =
    \sum_{k = 0}^{n} 3(k+1)q_{k + 1} x^{k}.
    $$
    Thus, for $Tq = 0,$ we need
    $$
    3(n + 1)q_{n + 1}x^{n} + \sum_{k = 0}^{n-1} [5(k+2)(k+1)q_{k + 2}
    - 3(k + 1)q_{k + 1}]  x^k = 0.
    $$
    Since $\mathcal P_{n}(\mathbb R)$ has a unique $0,$ we know
    $$
    q_{n + 1} = 0,
    $$
    and back-substituting yields $q_k = 0$ for all $1 \leq k \leq n,$
    meaning $\text{null}\ T = \{p: p = \mu, \mu \in \mathbb R\},$
    with dimension $0.$ Hence, by the fundamental theorem of linear maps,
    $$
    n + 2=\dim \mathcal P_{n + 1}(\mathbb R) = 1 + \dim\text{range}(T),
    $$
    meaning $$
    \dim \text{range}(T) = n + 1 = \dim \mathbb P_n(\mathbb R).
    $$
    Since the dimension of the range of $T$ is equal to its codomain,
    $T$ must be surjective. Thus, for all $p \in \mathcal
    P_{n}(\mathbb R),$ there is a $q \in \mathcal P_{n + 1}(\mathbb
    R) \subset \mathcal P(\mathbb R)$ such that $5q'' + 3q' = p.$
    Since any $p$ has some finite degree $m,$ any polynomial $p \in
    \mathcal P(\mathbb R)$ has some $q \in \mathcal P(\mathbb R)$
    such that $5q'' + 3q' = p.$  \qed

  \item Assume that we have $T \in \mathcal L(V, W)$ for finite
    dimensional vector spaces $V, W$ and that for some choice of
    basis for $V, W,$ the matrix $\mathcal M(T)$ has less than $\dim
    \text{range}(T)$ non-zero entries.

    Then, $\mathcal M(T)$ must have less than $\dim \text{range}(T)$
    non-zero columns, meaning the dimension of $\dim
    \text{col}(\mathcal M(T)) < \dim \text{range}(T),$ since none of
    the bases can be $0.$

    However, since $\dim \text{range}(T) = \dim \text{col}(\mathcal
    M(T)),$ this is a contradiction. \qed

  \item Assume we have some finite dimensional vector space $V,$ with
    some $T \in \mathcal L(V).$ Additionally, we have $$
    B = \mathcal M(T, (v_1, \ldots, v_n), (v_1, \ldots, v_n))
    $$
    and $C = \mathcal M(I, (u_1, \ldots, u_n), (v_1, \ldots, v_n)).$

    Let $$
    A = \mathcal M(T, (u_1, \ldots, u_n), (u_1, \ldots, u_n)).
    $$

    By definition of the identity matrix, we know $$
    u_j = \sum_{i = 1}^{n} C_{ij} v_i.
    $$

    Let us consider what the matrices $CA$ and $BC$ do to $e_k.$  By
    definition, also know
    $$
    T(v_j) = \sum_{i = 1}^{n} B_{i j}v_i, \qquad T(u_j) = \sum_{i =
    1}^{n} A_{ij} u_i.
    $$

    Thus, using our expression for $u_j,$ we get
    \begin{align*}
      T\left(\sum_{i = 1}^{n} C_{ij}v_i\right) &= \sum_{i = 1}^{n}
      \sum_{k = 1}^{n} A_{ij} C_{ki} v_{k} \\
      &= \sum_{i = 1}^{n}\sum_{k = 1}^{n} C_{ij}  B_{ki} v_k.
    \end{align*}
    Since $v_1, \ldots, v_n$ is a basis, we can equate coefficients, yielding
    $$
    \sum_{i = 1}^{n} A_{ij} C_{ki} = \sum_{i = 1}^{n} C_{ij} B_{ki}
    $$
    for all $k, j.$ Notice that the left hand side is exactly $C_{k
    \cdot}A_{\cdot j},$ and the right hand side is $B_{k \cdot } C_{\cdot j},$
    meaning
    $$
    (CA)_{kj} = C_{k \cdot } A_{ \cdot j} = B_{k \cdot } C_{ \cdot j}
    = (BC)_{kj}
    $$
    for all $k, j.$ Equivalently, the matrices $CA$ and $BC$ are
    identical. Since $C$ is invertible (identity map), we get
    $$
    CA = BC \Rightarrow A = C^{-1} BC
    $$
    as required. \qed
\end{enumerate}

\subsection{Evaluating Proofs}
This proof is incorrect as stated, since to construct $T_i$ we need
$u \neq 0$ otherwise no such $T_i$ exists, and $u \neq 0$ exists only
if $V \neq \{0\}.$

\subsection{Using Computers}
\begin{center}
  \includegraphics[width=0.9\linewidth]{Screenshot 2026-10-09 at
  12.16.39 pm.png}
\end{center}
\end{document}
